%%%%%%%%%%%%%%%%%%%%%%%%%%%%%%%%
% Chapter 2. Background
%%%%%%%%%%%%%%%%%%%%%%%%%%%%%%%%

\section{Background}\label{sec:background}

% ----------------------------------------------------------------
\subsection{Soft Prompt Methods for LLM Reasoning}\label{sec:rw_soft_prompt}
% ----------------------------------------------------------------

% 나중에 이쪽 내용 Introduction과 비교해서 줄일 수 있을 듯 (중복되는게 많아 보임)

\textit{Soft prompts} emerged as a parameter-efficient alternative to full fine-tuning. Prefix-Tuning \citep{li2021prefix} introduced the paradigm by prepending trainable continuous vectors to every transformer layer, Prompt Tuning \citep{lester2021power} simplified this to input-layer-only tuning and showed that it matches full fine-tuning at sufficient scale, and P-tuning \citep{liu2021gpt} extended the approach with an LSTM-based prompt encoder to model inter-position dependencies.

More recently, \textit{soft prompts} have been actively applied to LLM reasoning. TTSV \citep{kang2025model} optimizes prefix embeddings at test time via entropy minimization, steering the model toward high-certainty states to activate latent reasoning capabilities. SoftCoT \citep{xu2025softcot} replaces explicit chain-of-thought tokens with soft tokens generated by an assistant model, leveraging the stronger representational capacity of continuous representations over discrete language. LTPO \citep{ye2025thinking} optimizes latent thought embeddings per test instance through policy gradient with a confidence-based intrinsic reward. COCONUT \citep{hao2024training} feeds the model's own hidden states back as input embeddings, enabling reasoning in a continuous latent space. MemGen \citep{zhang2025memgen} demonstrates that embedding vectors can serve as experience memory by constructing latent token sequences that enrich the model's reasoning. Pause tokens \citep{goyal2024think} insert learnable tokens into the input to provide additional computation steps.

These approaches all inject out-of-vocabulary continuous embeddings optimized for task-specific objectives. Because evaluation centers on end-task accuracy, the contribution of the injection itself and that of the optimization are not separated, and the mechanisms by which such out-of-vocabulary continuous embeddings shape model behavior remain a separate and underexplored question.
% COCONUT offers deeper mechanistic analysis but is likewise trained for a specific continuous-reasoning objective, leaving the effect of injection itself --- decoupled from purpose-driven optimization --- unexamined.

% ----------------------------------------------------------------
\subsection{Noise Injection in Neural Networks}\label{sec:rw_perturbation}
% ----------------------------------------------------------------

Neural networks have incorporated noise at various stages. During training, dropout \citep{srivastava2014dropout} randomly deactivates hidden units to prevent overfitting, and NEFTune \citep{jain2024neftune} adds noise to token embeddings during instruction fine-tuning, reporting improvements in instruction following. These methods use noise as a regularizer and do not apply perturbations at inference time.

At inference time, perturbation-based approaches primarily target robustness and safety. Randomized smoothing \citep{cohen2019certified} injects Gaussian noise into inputs to obtain certified robustness guarantees. In the LLM setting, SmoothLLM \citep{robey2023smoothllm} perturbs input prompts at the character level to defend against adversarial jailbreaking. These methods require multiple noisy forward passes with output aggregation, and target prediction stability rather than output diversity.

In reinforcement learning, NoisyNet \citep{fortunato2018noisy} injects learnable noise into network weights to improve exploration, demonstrating that the location and structure of noise injection affect exploration quality. This motivation is closest to RSPs, which similarly target diversity, but differ in injection target (input embeddings vs.\ network weights) and require no learned noise parameters.

RSPs differ from these approaches in the following respects. \textbf{(1)} RSPs operate solely at inference without any training, whereas dropout and NEFTune inject noise during training, modifying the resulting model parameters accordingly. \textbf{(2)} RSPs preserve the original input and append fresh embedding vectors in a single forward pass, whereas robustness methods corrupt the input and aggregate over multiple noisy passes. \textbf{(3)} RSPs use purely random vectors sampled from the pretrained embedding statistics, whereas NoisyNet learns noise parameters through optimization.



% ----------------------------------------------------------------
% \subsection{Exploration in RL-based LLM Training}\label{sec:rw_exploration}
% ----------------------------------------------------------------
%
% GRPO \citep{shao2024deepseekmath} is a reinforcement learning algorithm for LLMs that eliminates the critic model and estimates advantages through group-relative comparison. For each prompt, GRPO samples a group of $G$ responses and computes advantages by normalizing their rewards within the group. The policy is updated via a clipped surrogate objective (full formulation in Appendix~\ref{app:grpo}). Since the advantage is computed as a normalized group statistic, effective learning requires diverse outcomes within each group: if all $G$ responses receive the same reward, the standard deviation approaches zero and the learning signal vanishes.
%
% A central challenge in GRPO-based training is entropy collapse, where the policy entropy decreases rapidly as training progresses, leading to premature convergence and reduced output diversity \citep{yu2025dapo}. As the policy becomes increasingly deterministic, sampled responses within each group become nearly identical, diminishing the quality of the group-relative advantage signal.
%
% DAPO \citep{yu2025dapo} addresses entropy collapse through several modifications to GRPO. First, it removes the KL divergence penalty, which is unnecessary in long chain-of-thought reasoning where the policy naturally diverges from the reference model. Second, it introduces Clip-Higher, which decouples the clipping range into $\varepsilon_{\text{low}}$ and $\varepsilon_{\text{high}}$ and increases $\varepsilon_{\text{high}}$ to allow low-probability tokens to increase in probability more freely, promoting exploration. Third, it applies Dynamic Sampling, which filters out prompt groups where all responses are correct or all incorrect, retaining only groups with mixed outcomes that provide effective gradient signals.
%
% CE-GPPO \citep{su2025cegppo} takes a different approach by analyzing the role of clipped tokens in entropy dynamics. In standard PPO and GRPO, tokens whose importance sampling ratio falls outside the clipping interval $[1\!-\!\varepsilon, 1\!+\!\varepsilon]$ have their gradients discarded. CE-GPPO reintroduces these gradients in a controlled manner through coefficients $\beta_1$ and $\beta_2$, which scale the gradients from tokens outside the lower and upper clipping boundaries respectively. A larger $\beta_1$ amplifies gradients from positive-advantage low-probability tokens, promoting exploration and slowing entropy decrease, while a larger $\beta_2$ amplifies gradients from negative-advantage low-probability tokens, accelerating convergence.
%
% These methods share a common strategy: they modify the policy optimization objective to manage the exploration-exploitation trade-off during training. However, the trajectory sampling process itself remains based on temperature sampling from the current policy, which preserves token rankings and cannot expand the support of the output distribution. Our findings indicate that RSPs have the potential to address this limitation at the sampling level, perturbing the initial output distribution and generating more diverse trajectories when combined with temperature sampling.

